A star has an apparent magnitude $m_{\mathrm{U}} = 15.0$ in the $U$-band. The
$U$-band filter is ideal, i.e., it has perfect (100\%) transmission within the
band and is completely opaque (0\% transmission) outside the band. The filter
is centered at 360\,nm, and has a width of 80\,nm. It is assumed that the star
also has a flat energy spectrum with respect to frequency. The conversion
between magnitude, $m$, in any band and flux density, $f$, of a star in Jansky
($1\,\mathrm{Jy} = 1 \times 10^{-26}\,\mathrm{W\,Hz^{-1}\,m^{-2}}$) is given by
\[ f = 3631 \times 10^{-0.4 m}\,\mathrm{Jy} \]

\begin{description}
\item[(T8.1)] Approximately how many $U$-band photons, $N_0$, from this star
  will be incident normally on a $1\,\mathrm{m}^2$ area at the top of the
  Earth's atmosphere every second?
\end{description}

This star is being observed in the $U$-band using a ground based telescope,
whose primary mirror has a diameter of 2.0\,m. Atmospheric extinction in
$U$-band during the observation is 50\%. You may assume that the seeing is
diffraction limited. Average surface brightness of night sky in $U$-band was
measured to be 22.0\,mag/arcsec$^2$.

\begin{description}
\item[(T8.2)] What is the ratio, $R$, of number of photons received per second
  from the star to that received from the sky, when measured over a circular
  aperture of diameter $2''$?
\item[(T8.3)] In practice, only 20\% of $U$-band photons falling on the primary
  mirror are detected. How many photons, $N_{\mathrm{t}}$, from the star are
  detected per second?
\end{description}
